\section{Analysis}
\label{sec:analysis}

\paragraph{Computational bottlenecks.}
\cref{tab:profile} breaks down one re-plan; the simulator is not the bottleneck (32\,ms per step against 814\,ms for the policy).
The vision encoder takes 230\,ms in the released bfloat16 weights but 39\,ms in float16: Turing GPUs have no bfloat16 tensor cores and emulate it.
The action expert accounts for 18\% of the FLOPs of a re-plan but 63\% of its latency ($10\times51$\,ms), and its step time is the same in both precisions: a pass over 50 action tokens is only 11.5\,GFLOP, and replaying its kernels from a CUDA graph cuts it from 53 to 9.4\,ms (\cref{sec:improve}): at batch one the GPU waits on kernel launches rather than on arithmetic.

\begin{table}[t]
\centering\footnotesize
\setlength{\tabcolsep}{3pt}
\begin{tabular}{@{}lrrr@{}}
\toprule
Component (one re-plan) & GFLOP & bf16 (ms) & fp16 (ms) \\
\midrule
Vision encoder, 2 cameras & 431 & 230 & 39 \\
Prefix incl.\ vision (32 LM layers) & 527 & 304 & 82 \\
One expert pass (flow step) & 11.5 & 51 & 52 \\
Re-plan, $S{=}10$ (released) & 642 & 814 & 605 \\
Simulator step (CPU + EGL) & -- & \multicolumn{2}{c}{32} \\
\bottomrule
\end{tabular}
\caption{Cost of one re-plan on a Tesla T4 (isolated components, CUDA events, median of 60 calls; \cref{tab:latency}, from a separate run, agrees within 4\% for a full re-plan).}
\label{tab:profile}
\end{table}

\paragraph{Successful and failed trajectories.}
Successes finish in a median of 100 steps; all 46 failures run to the 280-step limit after at least one gripper closure.
From the executed gripper commands (object poses are not logged; segments of at most 3 steps merged), 9 failures close the gripper three or more times, 9 others keep it closed at the limit, and 28 reopen it after one or two closures; task~3 (bowl on the cookie box) fails in 10 of 20 episodes.
Fresh flow noise at every re-plan makes consecutive chunks differ about as much as independent samples (first normalised action: 1.31 from the shifted previous chunk, 1.37 between two noises; 5 LIBERO-Object episodes).
The gripper command chatters accordingly: failed episodes contain 12.6 gripper segments of at most 3 steps per 100 steps, successful ones 2.7.

\paragraph{Strengths.}
\textbf{S1 Small footprint:} 1.16\,GB of weights and 21\,MB of activations per re-plan (float16) fit a 16\,GB T4 several times.
\textbf{S2 Prefix caching:} the expert reuses the cached prefix, so a flow step costs 52\,ms against 304\,ms for the prefix.
\textbf{S3 Few-step flow:} with $S{=}3$ the first action deviates from $S{=}10$ by 0.32, a quarter of the sampling spread (1.37), and SR holds down to $S{=}1$ (\cref{sec:improve}).

\paragraph{Limitations.}
\textbf{L1 Precision mismatch:} bfloat16 weights run emulated on Turing.
\textbf{L2 Launch-bound expert:} the flow loop dominates latency at batch one while using little arithmetic.
\textbf{L3 Need for fresh images:} executing 5 actions per chunk, or re-planning with fresh proprioception but an image up to 4 steps old, loses 12.0--12.5 points (\cref{sec:improve}), so an episode needs a full re-plan at every step.
\textbf{L4 A fresh sample at every step:} independent noise at every re-plan produces gripper chatter, concentrated in failures, although removing it alone does not raise SR measurably (\cref{sec:improve}).
